\chapter{Introducción}
\label{cap:introduccion}
Uno de los aspectos más interesantes del mundo de la informática es la evolución
de los dispositivos de entrada y salida, en particular aquellos dedicados a la
interacción del usuario con el ordenador.  Actualmente, la interacción se
realiza principalmente a través de teclado y ratón, aunque en el ámbito de
los videojuegos existen otros dispositivos de entrada como \textit{gamepads},
\textit{joysticks}, \textit{dance pads}, sistemas de volantes que simulan la
conducción de un coche, etc. Con estos dispositivos se puede interactuar con los
sistemas de una forma más natural, divertida e intuitiva.

Por supuesto, crear experiencias de usuario que utilicen estos dispositivos no
es trivial. Existen muchos sistemas y soluciones para utilizar los dispositivos
convencionales que son el teclado y el ratón, pero no ocurre lo mismo con los
dispositivos de entrada más especializados.

Este trabajo de fin de grado tiene como objetivo principal el desarrollo de un videojuego que
utilice como medio de entrada un sistema de detección de posturas del cuerpo humano a través de una
cámara. Para ello, será necesario tanto el desarrollo de un método para detectar la postura del
cuerpo humano como la creación de un videojuego que utilice esta información para interactuar con el
usuario.

El objetivo de este documento es presentar el trabajo realizado, detallando tanto el proceso de
desarrollo del sistema, como las decisiones tomadas y los resultados obtenidos. Se comienza
detallando los objetivos del proyecto y situación actual del ámbito del problema (capítulos dos,
tres y cuatro). A continuación se describen las herramientas, normas y metodologías utilizadas,
así como las alternativas consideradas (capítulos cinco y ocho). También se presentan los requisitos
iniciales del sistema y las restricciones de alcance del proyecto (capítulos seis y siete). Seguido de
esto,  se detalla la solución propuesta y el proceso de desarrollo del sistema, junto a los riesgos
detectados durante y tras el desarrollo (capítulos nueve, diez y once). Finalmente, se presentan las
conclusiones del proyecto y las posibles líneas de mejora y evolución del sistema (capítulo doce).

Además, este documento se complementa con una serie de anexos que profundizan en algunos aspectos
más técnicos del desarrollo del sistema, como pueden ser la planificación del proyecto o la
especificación de seguridad y privacidad. A continuación, se presenta un breve resumen de
cada uno de los documentos adicionales:
\begin{itemize}
    \item \textbf{Anexo I, Especificaciones del sistema}: En este anexo se detallan los requisitos
        funcionales y no funcionales del sistema, así como las especificaciones de los módulos
        que lo componen.

    \item \textbf{Anexo II, Análisis y diseño del sistema}: En este anexo se presenta el análisis y
        diseño del sistema. Incluye los diagramas de clases de análisis y diseño, el diagrama de
        despliegue del sistema y la realización de los casos de uso, entre otros.

    \item \textbf{Anexo III, Estimación del tamaño y esfuerzo}: En este anexo se presenta la
        planificación del proyecto, incluyendo la estimación del tamaño y esfuerzo del sistema,
        incluyendo la metodología de planificación utilizada y las tareas realizadas.

    \item \textbf{Anexo IV, Plan de seguridad}: En este anexo se presenta el plan de seguridad del
        sistema. Incluye la identificación de los riesgos de seguridad de la aplicación y las
        medidas de seguridad implementadas para mitigarlos.
\end{itemize}

La estructura de la documentación es la aprobada por el Colegio de Ingenieros Informáticos adaptada
a las necesidades del proyecto, que se recomienda seguir a partir del curso 2024/2025. Dicha
estructura se puede encontrar en la página web del grado de Ingeniería Informática de la
Universidad de Salamanca\footnote{%
\url{https://diaweb.usal.es/diaweb/archivos/61771471Documentacio__769_nMemoriaTFG_v2.0.pdf}}.
